\chapter{Results}
\label{ch:results}

This chapter reports the outputs of each workstream for the Guadalajara
Metropolitan Area (\zmg). Section order follows the pipeline: the estimated demand
surface (\cref{sec:res-w1}), its calibration (\cref{sec:res-w2}), the coverage-gap
diagnosis and its supervised interpretation (\cref{sec:res-w3}), the
Node--Place--People prioritization (\cref{sec:res-w4}), the generated corridors
(\cref{sec:res-w6}), the audit of the existing network (\cref{sec:res-w7}), and the
validation suite (\cref{sec:res-w8}).

\section{Estimated transit-demand surface (W1)}
\label{sec:res-w1}

The Tier-1 pipeline produces a transit-demand estimate for all \num{1881} \ageb{}s,
totalling \num{8471576} daily transit trip-ends. Household vehicle ownership varies
widely across the metropolis: the mean vehicle rate is \num{0.577}, giving a mean
transit propensity of \num{0.423}. Because demand is weighted by
$(1-v_i)$ (\cref{eq:transit-demand}), high-car-ownership \ageb{}s contribute
proportionally less latent transit demand even where raw trip generation is high.
\Cref{fig:demand} maps the resulting surface, which concentrates along the dense
central and eastern corridors of the metropolis.

\begin{figure}[htbp]
  \centering
  \includegraphics[width=0.82\linewidth]{zmg_demand.pdf}
  \caption{Estimated transit-demand surface (\ageb trips per day, log scale). Demand
    concentrates in the dense, lower-car-ownership central and eastern zones.}
  \label{fig:demand}
\end{figure}

\section{Gravity-model calibration (W2)}
\label{sec:res-w2}

Calibrating the distance-decay parameter against the 2022 Origin--Destination survey
(\cref{eq:calibration}) selects $\beta^\star = \num{1.2005}$, a markedly weaker decay
than the initial $\beta = \num{2.0}$ prior---\zmg commuters travel farther than the
prior assumed. Crucially, the calibrated value \emph{improves} fit on the corrected
\num{1881}-\ageb universe on every metric (\cref{tab:calibration}), reversing the
provisional conclusion drawn on the earlier, mis-filtered \ageb table. The calibrated
$\beta$ is adopted for the demand surface and all downstream workstreams.

\begin{table}[htbp]
\centering
\caption{Gravity-model calibration against the EOD 2022 survey.}
\label{tab:calibration}
\begin{tabular}{lS[table-format=4.1]S[table-format=1.4]}
\toprule
Model & {RMSE (trips)} & {$R^2$} \\
\midrule
Prior, $\beta = 2.0$        & 5088.2 & {--} \\
Calibrated, $\beta = 1.2005$ & 4524.9 & 0.2498 \\
\bottomrule
\end{tabular}
\end{table}

\section{Coverage-gap diagnosis (W3)}
\label{sec:res-w3}

Cumulative-opportunities accessibility (\cref{eq:accessibility}) combined with the
demand surface yields the coverage-gap index (\cref{eq:gap}). Classifying \ageb{}s by
joint demand/accessibility quintiles gives \num{390} High-gap (\SI{20.7}{\percent}),
\num{1413} Medium-gap, and \num{78} Low-gap areas (\cref{tab:gap}). As
\cref{fig:coverage-gap} showed, the High-gap areas form a periphery around a
well-served (Low-gap) core.

\begin{table}[htbp]
\centering
\caption{Coverage-gap classification of the \num{1881} \zmg \ageb{}s.}
\label{tab:gap}
\begin{tabular}{lS[table-format=4.0]S[table-format=2.1]}
\toprule
Category & {\ageb{}s} & {Share (\%)} \\
\midrule
High-gap   & 390  & 20.7 \\
Medium-gap & 1413 & 75.1 \\
Low-gap    & 78   & 4.1 \\
\bottomrule
\end{tabular}
\end{table}

Training a supervised model on the binary High-gap target (\cref{sec:w3}), with all
supply-side variables excluded, recovers the gap structure well and without leakage
(\cref{tab:retrain}). SHAP attribution ranks population, social lag
(\code{pe\_rezago}), and the employment proxy as the dominant drivers: High-gap areas
are dense, high-need, employment-rich zones that the current network under-serves.

\begin{table}[htbp]
\centering
\caption{Supervised retrain on the coverage-gap target (test fold).}
\label{tab:retrain}
\begin{tabular}{lS[table-format=1.3]S[table-format=1.3]}
\toprule
Model & {PR-AUC} & {ROC-AUC} \\
\midrule
Random Forest & 0.877 & 0.962 \\
LightGBM      & 0.883 & 0.962 \\
\bottomrule
\end{tabular}
\end{table}

\section{Node--Place--People prioritization (W4)}
\label{sec:res-w4}

The CRITIC$\,\times\,$EWM procedure (\cref{eq:critic}) assigns the largest objective
weights to the employment proxy (\num{0.163}), service density (\num{0.145}), and
four-way intersection density (\num{0.097}); the full weight vector appears in
\cref{fig:weights}. Averaged over all \ageb{}s, the place score is
$\overline{\mathrm{NPP}} = \num{0.4595}$, the equity term
$\overline{\mathrm{Eq}} = \num{0.2274}$, and the composite
$\overline{F} = \num{0.4131}$ at $\alpha = \num{0.20}$; the low equity mean is
expected, since \zmg is predominantly low-marginalization. A sensitivity analysis
over $\alpha \in \{0.10, 0.20, 0.30\}$ confirms the prioritization is robust: the
ranking correlates with the $\alpha = \num{0.20}$ baseline at Spearman
$\ge \num{0.98}$. This aggregate-correlation robustness, however, understates how much
$\alpha$ moves specific AGEBs in and out of the top of the list: the top-50 overlaps
the $\alpha=0.20$ baseline's top-50 at only \num{36}--\num{37} of \num{50} AGEBs across
the two off-baseline values, the top-100 at \num{78}--\num{82} of \num{100}, and
quintile membership changes for \SI{14}{\percent}--\SI{16}{\percent} of all AGEBs. The
equity weight therefore shifts which specific \ageb{}s top the list --- a
decision-relevant amount for a planning agency reading off a ranked shortlist --- even
though it does not disturb the broad structure the aggregate correlation reports.

\section{Generated corridors (W6)}
\label{sec:res-w6}

The generator produces five candidate corridors, of which four are feasible
(\cref{tab:corridors}); the fifth (\code{W6\_G05}) is rejected on the
anchor-directness constraint (\num{1.93} > \num{1.8}). \Cref{fig:corridors} maps the
four feasible corridors over the coverage-gap surface, and \cref{fig:pareto} places
all candidates in objective space.

\begin{table}[htbp]
\centering
\caption{Generated corridor candidates. \code{W6\_G05} fails the anchor-directness
  constraint.}
\label{tab:corridors}
\begin{tabular}{l S[table-format=2.1] S[table-format=2.0] S[table-format=6.0] l c}
\toprule
Corridor & {km} & {\ageb{}s} & {Demand/day} & Mode & Feasible \\
\midrule
W6\_G00 & 7.3  & 18 & 66041  & BRT             & Yes \\
W6\_G01 & 23.0 & 27 & 96839  & Light rail      & Yes \\
W6\_G02 & 12.1 & 25 & 192357 & Light rail      & Yes \\
W6\_G03 & 2.4  & 5  & 35784  & BRT             & Yes \\
W6\_G05 & 5.4  & 14 & 80234  & Light rail      & {--} \\
\bottomrule
\end{tabular}
\end{table}

\begin{figure}[htbp]
  \centering
  \includegraphics[width=0.78\linewidth]{w6_pareto.pdf}
  \caption{Corridor candidates in $(f_2, f_1)$ objective space, colored by
    feasibility. Feasible corridors span a range of length/gain trade-offs.}
  \label{fig:pareto}
\end{figure}

On the three-axis merit tests introduced for validation (\cref{sec:res-w8}), \code{W6\_G02}
is the only feasible corridor that passes all three: \SI{56}{\percent} of its served
\ageb{}s are High-gap (versus the \SI{20.7}{\percent} metropolitan baseline), it is
non-redundant on the served-\ageb Jaccard measure (best Jaccard \num{0.18}), and its demand
per kilometre sits at the \num{73}rd percentile of the existing network. A second,
shape-based overlap measure computed in the same validation pass is less clean:
\code{W6\_G02}'s alignment passes within \SI{400}{\metre} of the existing premium
route MP-C03 for \SI{42}{\percent} of its sampled length---a materially higher figure
than the AGEB-set Jaccard overlap, and read here as read-across corroboration that a
planner-drawn alignment traces similar ground, at the cost of a caveat on how
\emph{unique} the corridor's path actually is; both figures should be reported
together rather than the Jaccard figure alone.

This merit-test result should not, however, be read in isolation from the framework's own
formal evaluation function. \Cref{tab:pareto-w5} reports the raw W5 objective values and
non-dominated (Pareto) rank for every candidate. Under strict Pareto dominance,
\code{W6\_G03} is the sole rank-1 (non-dominated) feasible corridor: it strictly dominates
\emph{every other candidate, feasible or not}, on all three objectives simultaneously
(higher demand gain $f_1$, shorter length, higher mean equity $f_3$).
\code{W6\_G02}---the corridor the merit tests single out---is dominated twice over,
by both \code{W6\_G00} and \code{W6\_G03}, and ranks \emph{last} (Pareto rank 3) among the
four feasible corridors. This is not an artifact of the specific composite-score weights
(\num{0.50}/\num{0.25}/\num{0.25}, \cref{eq:objectives}): a full sweep of \num{231}
alternative non-negative weight combinations on a \num{0.05} simplex grid finds
\code{W6\_G03} top-scoring under \emph{every} combination tested, which follows directly
from strict Pareto dominance and would hold for any weighting.

\begin{table}[htbp]
\centering
\caption{W5 objective values and Pareto rank for all five candidates (2026-09-09,
  recomputed from and cross-checked against \code{features.route\_candidates}).
  Rank 1 = non-dominated. $f_1$ and $f_3$ are maximized, $f_2$ is minimized.}
\label{tab:pareto-w5}
\begin{tabular}{l S[table-format=1.3] S[table-format=2.2] S[table-format=1.3] c c}
\toprule
Corridor & {$f_1$ (demand gain)} & {$f_2$ (km)} & {$f_3$ (equity)} & Pareto rank & Feasible \\
\midrule
W6\_G00 & 0.301 & 7.32  & 0.338 & 2 & Yes \\
W6\_G01 & 0.270 & 23.00 & 0.347 & 2 & Yes \\
W6\_G02 & 0.164 & 12.13 & 0.265 & 3 & Yes \\
W6\_G03 & 0.357 & 2.44  & 0.349 & 1 & Yes \\
W6\_G05 & 0.125 & 5.40  & 0.271 & 2 & No  \\
\bottomrule
\end{tabular}
\end{table}

The two evaluation lenses disagree because they measure different things.
The merit tests reward absolute demand, absolute High-gap share, and demand efficiency
relative to the \emph{existing network}; the W5 objective $f_1$ is a demand-weighted
\emph{average} unserved fraction over the corridor's own served \ageb{}s---a rate, not a
total. \code{W6\_G03}'s five anchors are apparently uniformly high-gap, so its average is
high despite carrying a fraction of \code{W6\_G02}'s absolute demand (\num{35784} versus
\num{192357} trips/day); \code{W6\_G02} serves a demand-weighted \emph{mix} of gap levels
across its \num{25} \ageb{}s, which pulls its average down even though it addresses far
more total unmet demand in absolute terms.

This is testable rather than merely asserted. Define $f_1^{\text{tot}} = \sum_{i\in S} D_i\,
\kappa\, u_i$, the same numerator as \cref{eq:objectives} without dividing by
$\sum_{i\in S} D_i$---an absolute, not per-unit, demand-weighted gain. Recomputing Pareto
dominance on $(f_1^{\text{tot}}, f_2, f_3)$ for the four feasible corridors (2026-09-09,
via the production \code{evaluate\_objective} code path, not a re-implementation) gives
$f_1^{\text{tot}} = \num{19873}$, \num{26100}, \num{31608}, \num{12791} for
\code{W6\_G00}--\code{W6\_G03} respectively, and \emph{all four become mutually
non-dominated}---\code{W6\_G03} no longer dominates the set. This is the more telling
result: under the rate formulation, \code{W6\_G03} wins regardless of objective weights
not because it is unambiguously best, but because ``small'' and ``short'' move together
for it, leaving no real trade-off for the other objectives to break; a genuinely
competing multi-objective problem should not produce one candidate that dominates
outright under every possible weighting, and the fact that it did here was itself a sign
that $f_1$'s rate formulation was mismeasuring the trade-off, not that \code{W6\_G03}
was secretly the strongest candidate.

Correcting for this reframes the choice rather than resolving it to a single winner:
\code{W6\_G02} and \code{W6\_G01} address substantially more absolute unmet demand per
kilometre (\num{2606} and \num{1135} demand-units/km) than \code{W6\_G03}
(\num{12791}/2.44 = \num{5242}/km is highest, but on a corridor an order of magnitude
smaller), so the real, previously undisclosed choice is between a smaller number of
larger, more consequential corridors and a highly-targeted but minor one---a genuine
scale-versus-targeting-purity trade-off for a planning agency to make explicitly, not an
error the thesis needs to resolve on the reader's behalf. \Cref{sec:disc-generative}
returns to this with the recommendation that follows from it.

The anchor-directness substitution for endpoint-detour (\cref{sec:w6}) was motivated
by \code{W6\_G02}'s specific geometry, so it is worth asking whether its effect is
confined to that one case. \Cref{tab:directness-comparison} reports both metrics for
the full candidate set: it is not. Endpoint-detour would pass only \code{W6\_G03} (a
near-collinear, two-anchor stub) and reject \code{W6\_G00}, \code{W6\_G01}, and
\code{W6\_G02} alike, each exceeding the \num{1.8} cap by a wide margin
(\num{2.84}, \num{2.02}, and \num{2.09} respectively). Anchor-directness passes all
three. The substitution therefore has a broad, systematic effect across the
multi-anchor candidates---rescuing three corridors, not narrowly tuning one---which
is a materially different, and more defensible, story than a single worked example
can tell on its own.

\begin{table}[htbp]
\centering
\caption{Anchor-directness versus endpoint-detour for the full candidate set
  (route\_km unchanged between metrics; only the denominator differs).
  Recomputed 2026-09-09 by re-running the generator against the live database;
  the anchor-directness column and route lengths reproduce
  \cref{tab:corridors} exactly.}
\label{tab:directness-comparison}
\begin{tabular}{l S[table-format=2.2] S[table-format=1.2] S[table-format=1.2] c c}
\toprule
Corridor & {km} & {Anchor-directness} & {Endpoint-detour} & Anchor pass? & Endpoint pass? \\
\midrule
W6\_G00 & 7.32  & 1.44 & 2.84 & Yes & No \\
W6\_G01 & 23.00 & 1.16 & 2.02 & Yes & No \\
W6\_G02 & 12.13 & 1.54 & 2.09 & Yes & No \\
W6\_G03 & 2.44  & 1.25 & 1.36 & Yes & Yes \\
W6\_G05 & 5.40  & 1.93 & 2.69 & No  & No \\
\bottomrule
\end{tabular}
\end{table}

\section{Existing-route audit (W7)}
\label{sec:res-w7}

Scoring all \num{247} existing SITEUR routes with the same multi-objective function
flags \num{229} of them (\cref{tab:audit}): \num{109} as Indirect, \num{80} as Low
demand, and \num{40} as Redundant, leaving only \num{18} unflagged.
\Cref{fig:w7audit} plots every route in objective space by flag. Each flagged route
carries a modification proposal (shortcut, merge, or retire), making the audit
directly comparable to the generated corridors.

\begin{table}[htbp]
\centering
\caption{Audit flags across the \num{247} existing SITEUR routes.}
\label{tab:audit}
\begin{tabular}{lS[table-format=3.0]}
\toprule
Flag & {Routes} \\
\midrule
Indirect   & 109 \\
Low demand & 80 \\
Redundant  & 40 \\
Unflagged  & 18 \\
\bottomrule
\end{tabular}
\end{table}

\begin{figure}[htbp]
  \centering
  \includegraphics[width=0.82\linewidth]{w7_audit.pdf}
  \caption{Existing-route audit: the \num{247} SITEUR routes in $(f_2, f_1)$ space,
    colored by audit flag.}
  \label{fig:w7audit}
\end{figure}

\section{Validation (W8)}
\label{sec:res-w8}

\paragraph{Hold-out backtest and benchmark.} Masking the premium tier (Mi Macro and
Mi Tren) and re-running the generator reproduces the masked routes with a mean
shape overlap of only \SI{15.0}{\percent}; masking individual rail lines yields
\SI{0}{\percent}. The benchmark against \num{33} existing premium routes gives a
mean overlap of \SI{10.5}{\percent} for \SI{44.9}{\kilo\metre} of generated
corridor. Low overlap is the expected and desired outcome: the generator identifies
genuinely under-served areas rather than replicating existing lines. Reconstruction
is, in any case, a weak proxy for corridor merit (\cref{ch:discussion}).

\paragraph{Before/after metrics.} Adding the feasible corridors to the existing
network improves every coverage metric (\cref{tab:beforeafter}): the \ageb coverage
rate rises by \num{1.1} percentage points, the accessibility Gini falls by
\num{0.0187}, and \num{47} previously-unserved \ageb{}s---home to \num{120648}
people---gain service.

\begin{table}[htbp]
\centering
\caption{Before/after coverage metrics for the feasible W6 corridors.}
\label{tab:beforeafter}
\begin{tabular}{lS[table-format=2.1]S[table-format=2.1]S[table-format=+1.4]}
\toprule
Metric & {Before} & {After} & {$\Delta$} \\
\midrule
Coverage rate (\%)  & 69.9   & 71.0   & +1.1 \\
Accessibility Gini  & 0.6333 & 0.6146 & -0.0187 \\
\bottomrule
\end{tabular}
\end{table}

\paragraph{Out-of-sample corroboration.} The strongest validation is external: the
2024 \textsc{gtfs} snapshot predates the December~2025 opening of SITEUR Line~4, so
the diagnostic treats the Line~4 corridor as unserved. Its \ageb{}s show
\num{1.6} times the metropolitan High-gap rate, confirming out-of-sample that the
coverage-gap index flags corridors that planners independently chose to build. This
is a single natural experiment ($n=1$), reported here as a bare ratio with no
confidence interval or significance test against the base rate of a model that
already flags \SI{20.7}{\percent} of the metropolis as High-gap; it should be read
as directional corroboration, not as a statistically quantified test, and further
out-of-sample cases (as future lines open, or via the multi-line backtests of
\cref{ch:transferability}) would strengthen the claim considerably. It corroborates
the \emph{diagnostic} layer; it does not prove the generator would have produced
Line~4, a distinction developed in \cref{ch:discussion}.
